% Generated draft: requires mathematical and source review.
\documentclass[11pt]{article}
\usepackage[margin=25mm]{geometry}
\usepackage{fontspec}
\setmainfont{Latin Modern Roman}
\usepackage{amsmath,amssymb,mathtools,unicode-math}
\setmathfont{Latin Modern Math}
\usepackage{enumitem,needspace,xurl,hyperref}
\setlength{\parindent}{0pt}
\setlength{\parskip}{5pt}
\setlength{\emergencystretch}{3em}
\newcommand{\N}{\mathbb{N}}
\newcommand{\Z}{\mathbb{Z}}
\newcommand{\Q}{\mathbb{Q}}
\newcommand{\R}{\mathbb{R}}
\newcommand{\C}{\mathbb{C}}
\newcommand{\F}{\mathbb{F}}
\newcommand{\A}{\mathbb{A}}
\newcommand{\PP}{\mathbb{P}}
\newcommand{\E}{\mathbb{E}}
\newcommand{\eps}{\varepsilon}
\newcommand{\dd}{\,\mathrm{d}}
\begin{document}
\section*{UnsolvedMath ID: 30004040}
\textbf{Source problem code:} catalogue entry
\subsection{Definitions and mathematical statement}
Use the conventional setting of finite simple undirected graphs; this domain is not explicitly specified in the supplied record and requires confirmation against the original source. Let \(\mathfrak G\) denote the class of such graphs. For \(G\in\mathfrak G\) and \(v\in V(G)\), local complementation at \(v\) toggles adjacency between every pair of distinct neighbors of \(v\), leaving all other adjacencies unchanged. A graph \(H\) is a vertex-minor of \(G\) if it is isomorphic to a graph obtained from \(G\) by a finite sequence of local complementations and vertex deletions, allowing an empty sequence.

A graph class \(\mathcal F\subseteq\mathfrak G\) is closed under isomorphism if every graph isomorphic to a member also belongs to \(\mathcal F\). It is closed under vertex-minors if every vertex-minor of each member belongs to \(\mathcal F\). Interpret \emph{proper} as \(\mathcal F\subsetneq\mathfrak G\). Define the clique number by
\[
\omega(G)=\max\bigl\{|S|:S\subseteq V(G),\ \text{every two distinct vertices of }S\text{ are adjacent}\bigr\}.
\]
In particular, \(\omega(G)=0\) for the graph with no vertices.

\textbf{Question.} For every proper graph class \(\mathcal F\) closed under isomorphism and vertex-minors, does there exist an algorithm \(A_{\mathcal F}\) and a polynomial \(p_{\mathcal F}\) such that, for every input graph \(G\in\mathcal F\), the algorithm outputs \(\omega(G)\) within \(p_{\mathcal F}(|V(G)|+1)\) steps, using a standard explicit graph encoding?

Here \(\mathcal F\) is fixed before the algorithm is chosen, and both the algorithm and its polynomial bound may depend on \(\mathcal F\). Membership in \(\mathcal F\) is treated as an input promise; no recognition algorithm for \(\mathcal F\) is requested. These complexity conventions make the usual reading explicit, but the supplied record does not expressly discuss uniformity, encoding, or behavior on inputs outside \(\mathcal F\). The requested output is the size of a maximum clique. The supplied record labels the problem open; that label is not an independently verified mathematical conclusion.
\subsection{Short English statement}
For each fixed proper class of finite simple undirected graphs closed under vertex-minors and isomorphism, can the size of a maximum clique be computed in polynomial time on graphs promised to belong to the class? The graph domain and algorithmic conventions are standard interpretations requiring source confirmation.
\subsection{Sources}
\begin{itemize}
\item Supplied local UnsolvedMath registry record, permanent ID \(30004040\), title \emph{Maximum Clique in Proper Vertex-Minor-Closed Classes}, exact source problem code \texttt{catalogue entry}, in the Oberwolfach Reports Open Problems collection. No individual proposer is supplied; the recorded proposed year is 2019.
\item Supplied report citation: \emph{Graph Theory (2020)}, Oberwolfach Reports, DOI \texttt{10.4171/owr/2019/1}, pp.\ 5--63. \url{https://doi.org/10.4171/owr/2019/1}. The background describes Oberwolfach Reports as documenting mathematical workshops and their open-problem discussions. The citation's 2020 title date and the record's 2019 year information are retained without reconciliation.
\item Supplied literature reference: McCarty, \emph{Local Structure for Vertex-Minors}. \url{https://dspacemainprd01.lib.uwaterloo.ca/server/api/core/bitstreams/a3790341-2224-4c60-b8a6-34dc44a88c07/content}.
\item Supplied literature reference: \emph{Vertex-minors survey}. \url{https://sangiloum.github.io/legacy/2024/03/2023vertexminors-survey-revised.pdf}.
\item Literature assessment dated exactly \texttt{2026-08-22}: the supplied curation labels the problem open and states that the polynomial-time maximum-clique conjecture remains open beyond special classes such as circle graphs and bounded rank-width. This is supplied dated curation derived from the OWR material, not independently verified here. The background expressly requests independent verification of time-sensitive status and cited claims.
\item The registry assigns the statement status \texttt{corrected\_verified}, explaining that notation, typography, duplicate clauses, or omitted definitions were repaired while preserving mathematical content. This is a supplied editorial assertion, not independent mathematical or source verification of this draft.
\end{itemize}
\end{document}
